\chapter{Fisherの線形判別分析}
    \section{クラス間分散とクラス内分散の比を最大にする\texorpdfstring{$\bm{w}$}{w}の導出}
        \newcommand{\xbar}{\overline{\bm{x}}}
        \newcommand{\SW}{{S_\textrm{W}}}
        \newcommand{\SB}{S_\textrm{B}}
        クラス内共分散行列とクラス間共分散行列をそれぞれ$\SW,\;\SB \coloneqq (\xbar^1 - \xbar^2)(\xbar^1 - \xbar^2)^\top$とする。
        ここに$\xbar^1,\xbar^2$はクラス1,2のサンプルデータの平均である。
        $\SW$の正則性と$\xbar^1 \neq \xbar^2$を仮定すると、$\SW$は正定値対称、$\SB$は半正定値対称で階数は1となる。
        \par
        クラス間分散とクラス内分散の比
        \[ r(\bm{w}) \coloneqq \frac{\bm{w}^\top\SB\bm{w}}{\bm{w}^\top\SW\bm{w}} \]
        は$\bm{w}$の定数倍の差に影響されない。ある$\bm{w}^*$が$r$を最大化するならば任意の$\alpha\neq 0$に対して$\alpha\bm{w}^*$も$r$を最大化する。
        よって$\bm{w}^\top\SW\bm{w} = 1$という制約を加えても$r$の最大値は変わらないから、結局次の問題の解集合を含む最小の部分空間から$\bm{0}$を抜いたものが元の問題の解集合となる。
        \begin{align*}
            \textrm{maximize} &\quad f(\bm{w}) \coloneqq \bm{w}^\top\SB\bm{w} \\
            \textrm{subject to} &\quad g(\bm{w}) \coloneqq \bm{w}^\top\SW\bm{w} - 1 = 0
        \end{align*}
        実行可能領域は有界閉集合($n$次元の楕円)であり、目的関数は連続であるからこの最大化問題は解をもつ。
        $\SW$は正定値であるから実行可能領域に特異点は存在せず、Lagrangeの未定乗数法が使える。
        制約条件に対するLagrange定数を$\lambda$とすると、最適解の満たすべき必要条件は次のようになる。
        \begin{empheq}[left=\empheqlbrace]{align}
            (\SB - \lambda\SW)\bm{w} &= \bm{0} \\
            \bm{w}^\top\SW\bm{w} &= 1
        \end{empheq}
        (2)より$\bm{w}\neq\bm{0}$であるから、(1)の行列$\SB - \lambda\SW$が非正則になる必要がある。
        ここで$\lambda = 0$としてみると(1)より$\SB\bm{w} = \bm{0}$つまり$f(\bm{w}) = 0$となるが、$\SB$は階数1の半正定値行列なのでこのような$\bm{w}$は明らかに最大化問題の解ではない。
        ($\SB$の固有値の一つが正で他の$n-1$個が0なので$f(\bm{w})>0$となる$\bm{w}$が実行可能領域に必ず存在する。)
        よって$\lambda\neq 0$である。
        (1)の両辺に左から$\SW^{-1}$を掛けると上の連立方程式は次のようになる。
        \begin{empheq}[left=\empheqlbrace]{align}
            (\SW^{-1}\SB - \lambda I_n)\bm{w} &= \bm{0} \\
            \bm{w}^\top\SW\bm{w} &= 1
        \end{empheq}
        (3)より$\SW^{-1}\SB\bm{w} = \lambda\bm{w}$となるので$\lambda$は$\SW^{-1}\SB$の固有値であり、$\bm{w}$は対応する固有ベクトルである。
        \par
        突然だがここで
        \begin{align*}
            \bm{w}^* &= \frac{\SW^{-1}(\xbar^1 - \xbar^2)}{\sqrt{\left(\SW^{-1}(\xbar^1 - \xbar^2)\right)^\top\SW\left(\SW^{-1}(\xbar^1 - \xbar^2)\right)}} \quad (\neq\bm{0}) \\
            \lambda^* &= \left(\SW^{-1}(\xbar^1 - \xbar^2)\right)^\top\SW\left(\SW^{-1}(\xbar^1 - \xbar^2)\right) \quad (>0)
        \end{align*}
        とおいてみると
        \begin{align*}
            \SW^{-1}\SB\bm{w}^* &= \SW^{-1}(\xbar^1 - \xbar^2)(\xbar^1 - \xbar^2)^\top\frac{\SW^{-1}(\xbar^1 - \xbar^2)}{\sqrt{\left(\SW^{-1}(\xbar^1 - \xbar^2)\right)^\top\SW\left(\SW^{-1}(\xbar^1 - \xbar^2)\right)}} \\
            &= \sqrt{\left(\SW^{-1}(\xbar^1 - \xbar^2)\right)^\top\SW\left(\SW^{-1}(\xbar^1 - \xbar^2)\right)} \SW^{-1}(\xbar^1 - \xbar^2) \\
            &= \left(\SW^{-1}(\xbar^1 - \xbar^2)\right)^\top\SW\left(\SW^{-1}(\xbar^1 - \xbar^2)\right) \bm{w}^* \\
            &= \lambda^*\bm{w}^* \\
        \end{align*}
        となり、$\lambda^*, \bm{w}^*$はそれぞれ$\SW^{-1}\SB$の固有値,固有ベクトルになっている。
        さらに${\bm{w}^*}^\top\SW\bm{w}^* = 1$となる。
        よって力任せに作った$(\lambda^*, \bm{w}^*)$は(3),(4)の解である。
        同様に$(\lambda^*, -\bm{w}^*)$も解である。この2つ以外に解が存在しないことを示す。
        $\SW^{-1}\SB$は階数1であるから、Jordan分解すると次のようになる。
        \[
            \SW^{-1}\SB = JDJ^{-1},
            \quad D =
            \begin{bmatrix}
                d & 0 & \cdots & 0 \\
                0 & 0 & \cdots & \vdots \\
                \vdots & \vdots & \ddots & \vdots \\
                0 & \cdots & \cdots & 0
            \end{bmatrix},
            \quad d\neq 0
        \]
        ここに正則行列$J$はJordan標準形への適当な変換行列である。
        $\bm{z} = J^{-1}\bm{w}$すなわち$\bm{w} = J\bm{z}$なる変数変換を行うと次のようになる。
        \begin{align*}
            \left\{
                \begin{matrix}
                    (3) \\
                    (4)
                \end{matrix}
            \right. &\iff
            \left\{
                \begin{matrix}
                    (\SW^{-1}\SB - \lambda I_n)J\bm{z} &= \bm{0} \\
                    \bm{z}^\top J^\top\SW J\bm{z} &= 1
                \end{matrix}
            \right. \iff
            \left\{
                \begin{matrix}
                    J^{-1}(\SW^{-1}\SB - \lambda I_n)J\bm{z} &= \bm{0} \\
                    \bm{z}^\top J^\top\SW J\bm{z} &= 1
                \end{matrix}
            \right.
        \end{align*}
        \begin{empheq}[left = \iff\empheqlbrace]{align}
            (D - \lambda I_n)\bm{z} &= \bm{0} \\
            \bm{z}^\top J^\top\SW J\bm{z} &= 1
        \end{empheq}
        $\lambda \neq 0,d$のとき$D - \lambda I_n$は正則なので連立方程式は解を持たない。
        $\lambda = 0$は先述の通り最大化問題の解ではないので除外する。
        $\lambda = d$のとき$D - \lambda I_n$の階数は$n-1$となる。
        すなわち$D$の固有値$d$に対応する固有空間の次元は1であるから連立方程式の解は符号の違いによる2つしかない。
        よって先程作った$\lambda^* > 0$が実は$d$だったのであり、$(\lambda^*, \bm{w}^*)$と$(\lambda^*, -\bm{w}^*)$が連立方程式(3),(4)の解の全てである。
        \par
        Lagrangeの未定乗数法に基づく連立方程式(3),(4)は最大化問題の解の必要条件に過ぎないが、今回は解候補が2つしかなく、両方が同じ目的関数値を与えるのでこれが最適解である。
        \par
        ここまでは考察のために$\bm{w}^\top\SW\bm{w} = 1$という制約を付けていた。
        しかし冒頭で述べたように元の最大化問題では$\bm{w}$の定数倍の差は関係ないので、上の議論の$\bm{w}^*$ように規格化の手間を掛ける必要はなく、$\bm{w}$のノルムに関心が無いときは$\bm{w} = \SW^{-1}(\xbar^1 - \xbar^2)$とするだけで構わない。